\documentclass[12pt,a4paper]{article}
\usepackage[latin1]{inputenc}
\usepackage[spanish]{babel}
\usepackage{graphicx,amssymb,latexsym,amsmath,verbatim}
\usepackage{sans}
\usepackage{graphicx,amssymb,latexsym,amsmath}
\renewcommand{\thepage}{}
\renewcommand{\baselinestretch}{1}
\setlength{\parindent}{2em} \setlength{\textwidth}{18cm}
\setlength{\textheight}{24cm} \setlength{\topmargin}{-1cm}
\setlength{\oddsidemargin}{-1cm}
%\hyphenation{va-rie-da-des
%va-rie-dad pa-ra-le-lo a-si-mi-la a-pro-ve-cha-bles u-ni-for-me
%u-ni-for-me-men-te si-guien-te si-guien-tes re-su-mien-do
%ha-bi-tua-les}
%%%%%%%%%%%%%%%%%%%%%%%%%%%%%%%%%%%%%%%%%%%%%%%%%
%% Para empezar una interrogativa: ?`  %%%%%%%%%%
%%%%%%%%%%%%%%%%%%%%%%%%%%%%%%%%%%%%%%%%%%%%%%%%%
%% Para partir una matriz en cajas:         %%%%%
%% \[ \left(\begin{array}{c|c} A&B          %%%%%
%% \\ \cline{1-2} C&D\end{array}\right) \]  %%%%%
%%%%%%%%%%%%%%%%%%%%%%%%%%%%%%%%%%%%%%%%%%%%%%%%%


%SETS OF NUMBERS
%\newcommand{\R}{\mbox{${\bf R}$}}                  %reals
\newcommand{\R}{\mbox{${\mathbb R}$}}               %reals
%\newcommand{\C}{\mbox{${\bf C}$}}                  %complex
\newcommand{\C}{\mbox{${\mathbb C}$}}
%\newcommand{\N}{\mbox{${\bf N}$}}                  %natural numbers
\newcommand{\N}{\mbox{${\mathbb N}$}}
%\newcommand{\Z}{\mbox{${\bf Z}$}}                  %integers
\newcommand{\Z}{\mbox{${\mathbb Z}$}}
%\newcommand{\Q}{\mbox{${\bf Q}$}}                  %rationals
\newcommand{\Q}{\mbox{${\mathbb Q}$}}
\newcommand{\gp}{\mbox{${\bf F}_{p}$}}             % Fp
\newcommand{\pea}{\mbox{$ {\bf Z}_{p}$}}           %p-adic integers
\newcommand{\zn}{\mbox{${\bf Z}/{n}$}}             %Z/n
\newcommand{\zp}{\mbox{${\bf Z}/{p}$}}             %Z/p
\newcommand{\Zh}{\mbox{$\hat{{\bf Z}}$}}           %Z^
\newcommand{\xp}{\mbox{${\bf F}_{p}^{*}$}}         %units Fp

%GREEK LETTERS
\newcommand{{\ff}}{\mbox{$\varphi$}}                   %phi
\newcommand{\fp}{\mbox{$\varphi_{p}$}}               %phi_p
\newcommand{\al}{\mbox{$\alpha$}}                    %alpha
\newcommand{\be}{\mbox{$\beta$}}                     %beta
\newcommand{\px}{\mbox{$\varepsilon$}}               %epsilon
\newcommand{\Ga}{\mbox{$\Gamma$}}                    %capital gamma
\newcommand{\ga}{\mbox{$\gamma$}}                    %small gamma

%LOGIC SYMBOLS
\newcommand{\ex}{\mbox{$\exists$}}                   %there is
\newcommand{\fa}{\mbox{$\forall$}}                   %forall
\newcommand{\stf}{\mbox{$\models$}}                  %satisfies
\newcommand{\y}{\mbox{$\wedge$}}                     %and
%\newcommand{\o}{\mbox{$\vee$}}                       %or
\newcommand{\si}{\mbox{$\Longrightarrow$}}           %implies
\newcommand{\ekiv}{\mbox{$\Longleftrightarrow$}}     %equiv

%OTHERS
\newcommand{\dm}{\mbox{$\preceq$}}                   %dominates




\begin{document}

\small

\noindent{\sf Universidad Aut\'onoma de Madrid  \hfill �lgebra. Curso 2013--14}
\vspace{0.1cm}
\hrule
\vspace{0.1cm}
\noindent{\sf Ingenier�a Inform\'atica}
\vspace{0.1cm}
\begin{center}
{\bfseries  �LGEBRA. HOJA 7.}
\end{center}

\begin{enumerate}
%\renewcommand{\theenumi}{\bf\arabic{enumi})\ }
\renewcommand{\labelenumi}{\bf\arabic{enumi})\ }
%\renewcommand{\theenumii}{\alph{enumii})\ }
\renewcommand{\labelenumii}{\alph{enumii})\ }
\medskip




\item
Decide si el vector $\vec{v}=(1,0,1,0)$ es combinaci\'on lineal de los vectores $\{\vec{v}_1, \vec{v}_2, \vec{v}_3\} \subset \mathbb{R}^4$
para $$\vec{v}_1=(1,1,0,1),\ \vec{v}_2=(1,0,1,1)\ \ \hbox{y}\ \ \vec{v}_3=(1,0,0,-1).$$

 
\    
    
\item  Halla una base y calcula la dimensi�n de los siguientes subespacios vectoriales.

a)  $\mathbb S\subset \R^4$
definido por:
$$ \mathbb S= \left\{ \begin{array}{l} 
x_1-x_2+x_3-x_4=0\\
x_1-x_2=0 \\
x_3-x_4=0\\
x_1-x_4=0
 \end{array}
\right.
$$


b)  $\mathbb S\subset \R^4$
definido por:
$$ \mathbb S= \left\{ \begin{array}{l} 
x_1-x_2+x_3-x_4=0\\
x_1-x_2-2x_3-x_4=0 \\
 \end{array}
\right.
$$

 
\

\item  Halla una base de los siguientes subespacios vectoriales:

\vspace{0.15cm}

a)  $\mathbb S\subset \R^4$
generado por los vectores $\{\vec{v}_1, \vec{v}_2, \vec{v}_3, \vec{v}_4\} \subset \mathbb{R}^4$
donde
$$\vec{v}_1=(1,-1,0,0),\ \vec{v}_2=(2,0,-1,-1),\ \vec{v}_3=(1,0,0,-1),\ 
\vec{v}_4=(1,-1, 1, -1).$$

\vspace{0.15cm}

b)  $\mathbb S\subset \R^3$
generado por $\{\vec{v}_1, \vec{v}_2, \vec{v}_3, \vec{v}_4\} \subset \mathbb{R}^3$
siendo 
$$\vec{v}_1=(1,-1,0), \ \vec{v}_2=(1,0,-1), \ \vec{v}_3=(2,-1,-1),\ \vec{v}_4=(-1,-1, 2).$$

%%%%%% de la hoja 9 de 2011

\item
Sea $U$ el subespacio de $\R^4$ generado por $(1, 1, -1, 2)$ y $(-1, 2, 3, 3)$.
Sea $V$ el subespacio de $\R^4$ generado por $(1,7,3,12)$ y $(-4,5,1,1)$.
Sea $W = U+V$.

\textbf{a)}  Halla bases de los subespacios $U$, $V$ y $W$. Indica cu\'al es la dimensi\'on de $U\cap V$.
%\color{black}

\textbf{b)} Halla un sistema ecuaciones de $W$ con el menor n\'umero de ecuaciones.

\textbf{c)} Demuestra que $(2,-3,5,7) \notin W$.

\textbf{d)} Demuestra que el vector $\vec w = (27,0,27,57)$ est\'a en $W$ y halla $\vec u \in U$ y $\vec v \in V$ tales que $\vec w = \vec u + \vec v$.
 

\item
Sea $U$ el subespacio de $\R^4$ generado por los vectores $(1, 4, -2, 3)$ y $(2, 3, -1, 1)$.
Sea $V$ el subespacio de ecuaciones
{\setlength\arraycolsep{.1em}
\[
\left. \begin{array}{rrrrrrrrr}
 x_1 &-& 2x_2 &+& x_3 &+& 3x_4 &= & 0\\ 
 2x_1 &-& 2x_2 &-& x_3 &+& x_4 &=& 0\\
 x_1 &+& 3x_2 &-& 4x_3 &+& x_4 &=& 0
\end{array}
\right\}
\]
}%
Sea $W$ el subespacio suma $U+V$.

\textbf{a)} Demuestra que $W$ es la suma directa de $U$ y $V$, es decir, demuestra que $W = U \oplus V$.

\textbf{b)} Demuestra que el vector $\vec w = (19,-72,90,-75)$ est\'a en $W$.
Halla $\vec u \in U$ y $\vec v \in V$ tales que $\vec w = \vec u + \vec v$.

\item Sea $a \in \R$.
Sea $U$ el subespacio vectorial de $\R^4$ dado por las ecuaciones 
{\setlength\arraycolsep{.1em}
\[
\left. \begin{array}{rrrrrrrrr}
 x_1 &+& 2x_2 &-& x_3 &+& 2x_4 &=& 0\\ 
 2x_1 &+& x_2 &+& x_3 &-& x_4 &=& 0\\
\end{array}
\right\}
\]
}%
Sea $V$ el subespacio vectorial de $\R^4$ dado por las ecuaciones 
{\setlength\arraycolsep{.1em}
\[
\left. \begin{array}{rrrrrrrrr}
 3x_1 &+& x_2 & &  &+& 2x_4 &=& 0\\ 
 -x_1 & & &+& x_3 &+& a x_4 &=& 0\\
\end{array}
\right\}
\]
}%
\textbf{a)} Halla los valores de $a$ para los que $U \cap V \neq \{\vec 0\}$.
Para esos valores, calcula una base de $U \cap V$ y un sistema de ecuaciones de $U \cap V$ con el menor n\'umero de ecuaciones.

\textbf{b)} Para $a=-4$, halla $\vec u \in U$ y $\vec v \in V$ tales que $(3,1,-2,2) = \vec u + \vec v$.

\item Sean $U$ y $V$ subespacios vectoriales de $\R^n$ tales que $U \subset V$.
Halla $U \cap V$ y $U + V$.





\end{enumerate}






\end{document}

